\section{Rapport synthétique}

\subsection{Cas clinique}
Patient de $79$ ans est hospitalisé pour décompensation cardiaque avec fraction d’éjection conservée à $66$\%, Heart Failure with preserved Ejection Fraction (HFpEF) à la suite de l’administration de la première dose d’Azacitidine (Vidaza). Azacitidine est un agent déméthylant de l’ADN prescrit dans le cadre d’une leucémie myéloïde aiguë.\\
%\begin{singlespace}
Antécédents médicaux pertinents :

\begin{itemize}[label=\textbullet]
    \item Diabète de type $2$ ;
    \item Syndrome métabolique ;
    \item Cardiopathie hypertensive avec dysfonction diastolique \todoH{FloraStyle}{red!30}{\ !!! faut savoir c'est quoi la classification !}{\mbox{franche de grade $2$}} ;
    \item Dermo-hypodermite au niveau prétibial gauche ;
    \item Fibrillation atriale ;
    \item Infarctus du myocarde 
    \item BPCO GOLD II
    \item Leucémie myéloïde aiguë
    \end{itemize}
%\end{singlespace}

%\begin{singlespace}
Antécédents chirurgicaux pertinents :

\begin{itemize}[label=\textbullet]
    \item Mise en place d’un Pacemaker DDD ;
    \end{itemize}
%\end{singlespace}

%\begin{singlespace}
Allergies et assuétudes :

\begin{itemize}[label=\textbullet]
    \item Pas d’allergie connue à ce jour ;
    \item Tabac : $50$ UPA ($2$ paquets/j pendant $25$ ans) arrêté depuis une vingtaine d’années ;
    \item Alcool : occasionnelle ;
    \item Il nie toute autre consommation ;
    \end{itemize}
%\end{singlespace}



%\begin{singlespace}

Médication (en septembre $2026$) reprise dans le~\cref{tab:medication_patient} :

% ====================================================================================
% VERSION 3 : PAYSAGE / PENCHÉ (STYLE LES_HYPOTHESES_DIAGNOSTIQUES)
% Page pivotée automatiquement dans le visualiseur PDF avec grille complète (| et \hline)
% ====================================================================================
\begin{landscape}
\begingroup
\footnotesize
\setlength{\tabcolsep}{5pt}
\renewcommand{\arraystretch}{1.05}

\begin{xltabular}{\linewidth}{|>{\raggedright\arraybackslash\bfseries}p{4.2cm}|>{\raggedright\arraybackslash}p{3.6cm}|>{\raggedright\arraybackslash}p{5.8cm}|>{\raggedright\arraybackslash}X|}
    \caption{Médication du patient (septembre 2026)} \label{tab:medication_patient} \\
    \hline
    \rowcolor{gray!15}
    \textbf{DCI} & \textbf{Nom commercial \& Posologie} & \textbf{Classe pharmacothérapeutique} & \textbf{Indication clinique} \\
    \hline
    \endfirsthead

    \caption*{Tableau~\ref{tab:medication_patient} (suite) -- Médication du patient (septembre 2026)} \\
    \hline
    \rowcolor{gray!15}
    \textbf{DCI} & \textbf{Nom commercial \& Posologie} & \textbf{Classe pharmacothérapeutique} & \textbf{Indication clinique} \\
    \hline
    \endhead

    \hline
    \multicolumn{4}{|r|}{\footnotesize\textit{Suite du tableau à la page suivante...}} \\
    \hline
    \endfoot

    \hline
    \endlastfoot

    Ézétimibe / Atorvastatine & 
    Atozet 10\,mg / 40\,mg \newline 1 cp/j & 
    Hypolipémiant combiné (inhibiteur sélectif de l'absorption du cholestérol + statine de haute intensité) & 
    Prévention cardiovasculaire secondaire. \\
    \hline

    Bumétanide & 
    Burinex 1\,mg \newline 1 cp/j & 
    Diurétique de l'anse de Henle & 
    Traitement déplétif de la surcharge hydrosodée : contrôle des pressions de remplissage ventriculaire gauche élevées liées à la cardiopathie hypertensive avec dysfonction diastolique de grade 2. \\
    \hline

    Bisoprolol & 
    Bisoprolol 10\,mg \newline ½ cp/j (5\,mg/j) & 
    Bêtabloquant $\beta_1$ cardiosélectif & 
    Cardiopathie hypertensive \& contrôle hémodynamique : baisse des résistances vasculaires périphériques, allongement du temps de remplissage diastolique par ralentissement de la fréquence cardiaque, et freinage de l'hyperactivité sympathique. \\
    \hline

    Amiodarone & 
    Cordarone 200\,mg \newline 1 cp/j & 
    Antiarythmique de classe III de Vaughan-Williams & 
    Prévention des récidives de tachyarythmie supraventriculaire. \\
    \hline

    Fluticasone furoate / Uméclidinium / Vilantérol & 
    Trelegy Ellipta 92/55/22\,\textmu g \newline 1 inh/j & 
    Trithérapie inhalée fixe : Corticoïde inhalé (CSI) + Anticholinergique à longue durée d'action (LAMA) + $\beta_2$-mimétique à longue durée d'action (LABA) & 
    Traitement de fond de la BPCO. \\
    \hline

    Bromure d'ipratropium / Fénotérol & 
    Duovent \newline (posologie à la demande) & 
    Bronchodilatateur d'action rapide combiné : Anticholinergique à courte durée d'action (SAMA) + $\beta_2$-mimétique à courte durée d'action (SABA) & 
    Traitement de secours / à la demande : levée immédiate du bronchospasme et prise en charge des accès dyspnéiques paroxystiques dans la BPCO. \\
    \hline

    Apixaban & 
    Eliquis 5\,mg \newline 1 cp 2$\times$/j & 
    Anticoagulant oral direct (AOD) & 
    Thromboprophylaxie systémique liée à l'arythmie supraventriculaire associée (score $\mathrm{CHA_2DS_2}\text{-}\mathrm{VA}$ élevé). \\
    \hline

    Cholécalciférol (Vitamine $\mathrm{D_3}$) & 
    Vista $\mathrm{D_3}$ \newline 1 cp/j & 
    Vitamine D / régulateur du métabolisme phosphocalcique & 
    Prévention de l'ostéopénie / ostéoporose. \\
    \hline

    Zolpidem & 
    Stilnox 10\,mg \newline 1 cp le soir & 
    Hypnotique apparenté aux benzodiazépines & 
    Traitement symptomatique de courte durée de l'insomnie d'endormissement. \\
    \hline

    Azacitidine & 
    Vidaza \newline Cure palliative par voie sous-cutanée : 7 jours/mois pendant 28 jours & 
    Agent déméthylant de l'ADN & 
    Traitement de la leucémie myéloïde aiguë non éligible à une chimiothérapie d'induction intensive. \\
    \hline

    Dapagliflozine & 
    Forxiga 10\,mg \newline 1 cp/j & 
    Inhibiteur du cotransporteur sodium-glucose de type 2 (iSGLT2) & 
    Triple indication : \newline
    \textbullet~Diabète de type 2 ; \newline
    \textbullet~Cardioprotection ; \newline
    \textbullet~Néphroprotection. \\
    \hline

    Oxygénothérapie chronique & 
    2\,L/min par oxyconcentrateur & 
    Support ventilatoire & 
    Support pour l'hypoxémie nocturne et les troubles de diffusion dans le cadre de la BPCO. \\
\end{xltabular}
\endgroup
\end{landscape}

%\end{singlespace}

%\begin{singlespace}
Contexte psyho social : 
Le patient vit au domicile avec son épouse. Elle est très présente pour le patient. Le patient a $2$ enfants et $4$ petits-enfants. 
Le patient se sent très limité, dyspnéique et trouve le temps long en chambre.  Une thymie triste domine dans le discours du patient. 

%\end{singlespace}
%\begin{singlespace}
Histoire clinique :

Le patient est hospitalisé dans le service d’onco-hématologie pour insuffisance respiratoire de type I sur décompensation cardiaque. Suite à l’administration de la première dose d’Azacitidine (Vidaza) dans le contexte de sa leucémie myéloïde aiguë, le patient a présenté une dyspnée de grade IV NYHA évolutive avec sensation d’oppression thoracique. Ceci s’ajoutant à des signes cliniques caractéristiques de la décompensation cardiaque décrits dans l’examen clinique. Il a été décrit dans la littérature une augmentation significative d’événement cardio-pulmonaire lorsque le patient a des polymorbidités de ces systèmes. De plus, la dyspnée est répertoriée comme un effet indésirable très fréquent selon CBIP 
%\end{singlespace}
%\begin{singlespace}
Paramètres d’entrée :
\begin{itemize}[label=\textbullet]
    \item Tension artérielle couché : $96/70$
    \item Saturation : $91$\% sous $4$L/min d’oxygène
    \item Température : $36,8$°C
    \item Pouls : $69$/min
    \item Taille : $167$cm
    \item Poids : $58$kg
    \item BMI : $20,8$ kg/m$^2$
    \end{itemize}
%\end{singlespace}
%\begin{singlespace}
Examen clinique d’entrée :
    \begin{itemize}[label=\textbullet]
    \item Apparence : pâle.
    \item Examen neurologique : le patient est orienté dans l’espace et le temps. Score de Glasgow de $15/15$.
    \item Examen pulmonaire : eupnéique sous 4L/min d’oxygène. A l’auscultation, murmure vésiculaire diffus et symétrique sans sibilance. Crépitants bibasilaires entendus.
    \item Examen cardiaque et vasculaire : les jugulaires sont turgescentes avec un reflux hépato-jugulaire. Présence d’œdème des membres inférieurs de façon bilatérale avec prise du Godet à la palpation. A l’auscultation, B$1$B$2$ réguliers. Pas de souffle ni galop.
    \item Examen abdominal : péristaltisme normal à l’auscultation. Abdomen souple, dépressible, légèrement sensible en fosse iliaque gauche à la palpation. Pas d’organomégalie palpée.
    \end{itemize}
%\end{singlespace}

\subsection{Prise en charge hospitalière}
%\begin{singlespace}
%\end{singlespace}